\clearpage
\chapter*{KẾT LUẬN VÀ HƯỚNG PHÁT TRIỂN}
\phantomsection
\addcontentsline{toc}{chapter}{KẾT LUẬN VÀ HƯỚNG PHÁT TRIỂN}

\section*{Kết luận}
\phantomsection
\addcontentsline{toc}{section}{Kết luận}

Báo cáo đã xác định bài toán quản lý vòng đời chỗ ở tại ký túc xá PTIT và giới hạn phạm vi thành mười chức năng, chia theo năm giai đoạn: xem phòng, đặt phòng, thuê phòng, đóng phí và trả phòng. Các quy tắc nghiệp vụ then chốt đã được hệ thống hóa (BR01--BR12), bảo đảm một giường không bị giữ trùng, một sinh viên không có hai hợp đồng còn hiệu lực và các thay đổi trạng thái được thực hiện đồng bộ.

Phần phân tích đã xây dựng danh sách Use Case, ma trận tác nhân -- Use Case và đặc tả chi tiết các chức năng UC03, UC04, UC05, UC06, UC07, UC08, UC09. Định hướng công nghệ được lựa chọn gồm ReactJS cho giao diện, Java cho backend, PostgreSQL cho dữ liệu chính và Redis cho hỗ trợ tra cứu, giữ chỗ tạm.

\section*{Hướng phát triển}
\phantomsection
\addcontentsline{toc}{section}{Hướng phát triển}

Trong các giai đoạn tiếp theo, nhóm sẽ hoàn thiện các mục còn thiếu ở Chương 2 (đặc tả UC01, UC02, UC10 và các sơ đồ lớp, biểu đồ Use Case), sau đó chuyển sang thiết kế kiến trúc, cơ sở dữ liệu, giao diện và cài đặt minh họa luồng dữ liệu từ tạo phòng đến hoàn tất trả phòng.

Các hướng mở rộng gồm tích hợp cổng thanh toán thực tế, thông báo tự động, quản lý điện nước, tài sản, sửa chữa cơ sở vật chất, vi phạm nội trú và đồng bộ với hệ thống đào tạo, trên cơ sở vẫn giữ nguyên các ràng buộc toàn vẹn dữ liệu đã phân tích.
